% GENERATED FILE. Edit canonical lessons, then run npm run generate:books.
% canonical-source-hash: fnv1a64:4810c923775f61c6
% canonical-lessons: BN-C433-jhil, BN-C433-nala, BN-C433-jolabhumi, BN-C433-dol-khaoya, BN-C433-kolkol-kora

\chapter{Lakes, Ditches and Wetlands}
\label{ch:bn-lakes-ditches-and-wetlands}
\hlchaptermodality{433}

\begin{chapteropening}
\textbf{By the end of this chapter:} \emph{I can describe a lake, a ditch and a wetland, and water moving and murmuring.}

Complete the last lesson of chapter 433: I can describe a lake, a ditch and a wetland, and water moving and murmuring.
\end{chapteropening}

\section[\texorpdfstring{\bn{ঝিল} (jhil)}{ঝিল (jhil)}]{\bn{ঝিল} (jhil) — a lake; a marshy pool}

\label{lesson:BN-C433-jhil}



\begin{quote}
Before the new one: say the Bengali for a desert, then the Bengali for a jungle.
\end{quote}

\subsection*{You'll want to know: \bn{ঝিল}}
\textbf{\bn{ঝিল}} — \emph{jhil} — \textquotedblleft{}a lake; a marshy pool\textquotedblright{}.

\begin{quote}
\bn{শীতে} \bn{ঝিলে} \bn{অনেক} \bn{পাখি} \bn{আসে।}
\end{quote}

\begin{quote}
\emph{shīte jhile ônek pākhi āse.}
\end{quote}

\begin{quote}
In winter many birds come to the lake.
\end{quote}

\bn{হ্রদ} is a lake in geography; \bn{ঝিল} is the pool or marsh lake near a village.

\begin{cousinweb}
\textbf{\bn{ঝিল}} (\emph{jhil}) is inherited: CDIAL lists Bengali \emph{jhil}, \textquotedblleft{}a lake, a pool\textquotedblright{}, under the reconstructed \emph{jhilla} (CDIAL 5392).
\end{cousinweb}

\begin{grammarlens}[title={What you've built}]
One more word for getting around.
\end{grammarlens}

\subsection*{Your turn}
\begin{itemize}
\raggedright
  \item \emph{Say it:} \emph{jhil}
  \item \emph{Say it:} \emph{jhil}, once more
  \item \emph{Say it:} \emph{shīte jhile ônek pākhi āse.} — In winter many birds come to the lake.
  \item \emph{Recall:} say the Bengali for a desert, then the Bengali for a jungle, then say \emph{jhil} again
\end{itemize}

\begin{tcolorbox}[breakable,colback=teal!4,colframe=teal!35!black,title={Before you move on}]
What does \bn{ঝিল} mean? (\textquotedblleft{}A lake; a marshy pool\textquotedblright{}.) Say it once more.
\end{tcolorbox}
\section[\texorpdfstring{\bn{নালা} (nālā)}{নালা (nālā)}]{\bn{নালা} (nālā) — a ditch; a small channel}

\label{lesson:BN-C433-nala}



\begin{quote}
Before the new one: say the Bengali for a jungle, then the Bengali for a lake (natural).
\end{quote}

\subsection*{You'll want to know: \bn{নালা}}
\textbf{\bn{নালা}} — \emph{nālā} — \textquotedblleft{}a ditch; a small channel\textquotedblright{}.

\begin{quote}
\bn{বৃষ্টির} \bn{পরে} \bn{নালায়} \bn{জল} \bn{ভরে} \bn{গেল।}
\end{quote}

\begin{quote}
\emph{brishṭir pôre nālāy jôl bhôre gelô.}
\end{quote}

\begin{quote}
After the rain the ditch filled with water.
\end{quote}

\begin{cousinweb}
\textbf{\bn{নালা}} (\emph{nālā}) is inherited: CDIAL lists Bengali \emph{nālā}, \textquotedblleft{}a ditch\textquotedblright{}, under Sanskrit \emph{nāḍ\'{\={\i}}}, \textquotedblleft{}a tube, a pipe\textquotedblright{} (CDIAL 7047).
\end{cousinweb}

\begin{grammarlens}[title={What you've built}]
One more word for getting around.
\end{grammarlens}

\subsection*{Your turn}
\begin{itemize}
\raggedright
  \item \emph{Say it:} \emph{nālā}
  \item \emph{Say it:} \emph{nālā}, once more
  \item \emph{Say it:} \emph{brishṭir pôre nālāy jôl bhôre gelô.} — After the rain the ditch filled with water.
  \item \emph{Recall:} say the Bengali for a jungle, then the Bengali for a lake (natural), then say \emph{nālā} again
\end{itemize}

\begin{tcolorbox}[breakable,colback=teal!4,colframe=teal!35!black,title={Before you move on}]
What does \bn{নালা} mean? (\textquotedblleft{}A ditch; a small channel\textquotedblright{}.) Say it once more.
\end{tcolorbox}
\section[\texorpdfstring{\bn{জলাভূমি} (jôlābhūmi)}{জলাভূমি (jôlābhūmi)}]{\bn{জলাভূমি} (jôlābhūmi) — a wetland; marshland}

\label{lesson:BN-C433-jolabhumi}



\begin{quote}
Before the new one: say the Bengali for a lake (natural), then the Bengali for a ditch.
\end{quote}

\subsection*{You'll want to know: \bn{জলাভূমি}}
\textbf{\bn{জলাভূমি}} — \emph{jôlābhūmi} — \textquotedblleft{}a wetland; marshland\textquotedblright{}.

\begin{quote}
\bn{জলাভূমিতে} \bn{অনেক} \bn{পাখি} \bn{বাসা} \bn{বাঁধে।}
\end{quote}

\begin{quote}
\emph{jôlābhūmite ônek pākhi bāsā bā̃dhe.}
\end{quote}

\begin{quote}
Many birds nest in the wetland.
\end{quote}

\begin{cousinweb}
\textbf{\bn{জলাভূমি}} (\emph{jôlābhūmi}) joins \bn{জলা}, which CDIAL lists as Bengali \emph{jalā}, \textquotedblleft{}a marsh\textquotedblright{}, beside \emph{jal}, water (CDIAL 5155), to \emph{bhūmi}, \textquotedblleft{}land\textquotedblright{} (Monier-Williams).
\end{cousinweb}

\begin{grammarlens}[title={What you've built}]
One more word for getting around.
\end{grammarlens}

\subsection*{Your turn}
\begin{itemize}
\raggedright
  \item \emph{Say it:} \emph{jôlābhūmi}
  \item \emph{Say it:} \emph{jôlābhūmi}, once more
  \item \emph{Say it:} \emph{jôlābhūmite ônek pākhi bāsā bā̃dhe.} — Many birds nest in the wetland.
  \item \emph{Recall:} say the Bengali for a lake (natural), then the Bengali for a ditch, then say \emph{jôlābhūmi} again
\end{itemize}

\begin{tcolorbox}[breakable,colback=teal!4,colframe=teal!35!black,title={Before you move on}]
What does \bn{জলাভূমি} mean? (\textquotedblleft{}A wetland; marshland\textquotedblright{}.) Say it once more.
\end{tcolorbox}
\section[\texorpdfstring{\bn{দোল} \bn{খাওয়া} (dol khāoyā)}{দোল খাওয়া (dol khāoyā)}]{\bn{দোল} \bn{খাওয়া} (dol khāoyā) — to sway; to rock; to swing}

\label{lesson:BN-C433-dol-khaoya}



\begin{quote}
Before the new one: say the Bengali for a ditch, then the Bengali for a wetland.
\end{quote}

\subsection*{You'll want to know: \bn{দোল} \bn{খাওয়া}}
\textbf{\bn{দোল} \bn{খাওয়া}} — \emph{dol khāoyā} — \textquotedblleft{}to sway; to rock; to swing\textquotedblright{}.

\begin{quote}
\bn{হাওয়ায়} \bn{ঝিলের} \bn{নৌকাটা} \bn{দোল} \bn{খেল।}
\end{quote}

\begin{quote}
\emph{hāoyāy jhiler noukāṭā dol khelô.}
\end{quote}

\begin{quote}
The boat on the lake rocked in the breeze.
\end{quote}

Past (he/she …): \bn{সেই} \bn{নৌকা} \bn{দোল} \bn{খেল}, \emph{sei noukā dol khelô}.

\bn{দোল} is a swing; \bn{খাওয়া} here is to undergo, as in \bn{ধাক্কা} \bn{খাওয়া}.

\begin{cousinweb}
\textbf{\bn{দোল} \bn{খাওয়া}} (\emph{dol khāoyā}) joins \bn{দোল}, \textquotedblleft{}a swing, a swaying\textquotedblright{}, from Sanskrit \emph{dolā}, \textquotedblleft{}a swing, a litter\textquotedblright{} (Monier-Williams), to \bn{খাওয়া}, to eat, here to undergo.
\end{cousinweb}

\begin{grammarlens}[title={What you've built}]
One more word for getting around.
\end{grammarlens}

\subsection*{Your turn}
\begin{itemize}
\raggedright
  \item \emph{Say it:} \emph{dol khāoyā}
  \item \emph{Say it:} \emph{dol khāoyā}, once more
  \item \emph{Say it:} \emph{hāoyāy jhiler noukāṭā dol khelô.} — The boat on the lake rocked in the breeze.
  \item \emph{Recall:} say the Bengali for a ditch, then the Bengali for a wetland, then say \emph{dol khāoyā} again
\end{itemize}

\begin{tcolorbox}[breakable,colback=teal!4,colframe=teal!35!black,title={Before you move on}]
What does \bn{দোল} \bn{খাওয়া} mean? (\textquotedblleft{}To sway; to rock; to swing\textquotedblright{}.) Say it once more.
\end{tcolorbox}
\section[\texorpdfstring{\bn{কলকল} \bn{করা} (kôlkôl kôrā)}{কলকল করা (kôlkôl kôrā)}]{\bn{কলকল} \bn{করা} (kôlkôl kôrā) — to babble; to murmur (of running water)}

\label{lesson:BN-C433-kolkol-kora}



\begin{quote}
Before the new one: say the Bengali for a wetland, then the Bengali for to sway.
\end{quote}

\subsection*{You'll want to know: \bn{কলকল} \bn{করা}}
\textbf{\bn{কলকল} \bn{করা}} — \emph{kôlkôl kôrā} — \textquotedblleft{}to babble; to murmur (of running water)\textquotedblright{}.

\begin{quote}
\bn{ঝরনার} \bn{জল} \bn{কলকল} \bn{করছিল।}
\end{quote}

\begin{quote}
\emph{jhôrnār jôl kôlkôl kôrchhilô.}
\end{quote}

\begin{quote}
The water of the waterfall was babbling.
\end{quote}

Past (he/she …): \bn{সেই} \bn{জল} \bn{কলকল} \bn{করল}, \emph{sei jôl kôlkôl kôrlô}.

\begin{cousinweb}
\textbf{\bn{কলকল} \bn{করা}} (\emph{kôlkôl kôrā}) joins \bn{কলকল}, an imitative word for the sound of running water, to \bn{করা}, to do. No dictionary we use records an older form, so its history is not settled.
\end{cousinweb}

\begin{grammarlens}[title={What you've built}]
One more word for getting around.
\end{grammarlens}

\subsection*{Your turn}
\begin{itemize}
\raggedright
  \item \emph{Say it:} \emph{kôlkôl kôrā}
  \item \emph{Say it:} \emph{kôlkôl kôrā}, once more
  \item \emph{Say it:} \emph{jhôrnār jôl kôlkôl kôrchhilô.} — The water of the waterfall was babbling.
  \item \emph{Recall:} say the Bengali for a wetland, then the Bengali for to sway, then say \emph{kôlkôl kôrā} again
\end{itemize}

\begin{tcolorbox}[breakable,colback=teal!4,colframe=teal!35!black,title={Before you move on}]
What does \bn{কলকল} \bn{করা} mean? (\textquotedblleft{}To babble; to murmur (of running water)\textquotedblright{}.) Say it once more.
\end{tcolorbox}
